\documentclass[11pt]{article}
\usepackage[margin=1in]{geometry}
\usepackage{amsmath,amssymb,amsthm,mathtools,longtable,array}
\usepackage[hidelinks]{hyperref}
\allowdisplaybreaks[2]
\setlength{\emergencystretch}{3em}
\newtheorem{theorem}{Theorem}[section]
\newtheorem{lemma}[theorem]{Lemma}
\newtheorem{proposition}[theorem]{Proposition}
\newtheorem{corollary}[theorem]{Corollary}
\theoremstyle{remark}\newtheorem{remark}[theorem]{Remark}
\newcommand{\E}{\mathbb E}
\newcommand{\dd}{\,d}
\newcommand{\KK}{\mathcal K}
\newcommand{\BB}{\mathcal B}
\newcommand{\Gp}{G_+}
\newcommand{\eps}{\varepsilon}
\newcommand{\one}{\mathbf1}
\title{Full fixed-chart c-form assembly\\and a conditional CLOCK passage}
\author{Analytical increment for review}
\date{Version 1}
\begin{document}\maketitle
\begin{abstract}
On S2CP's stopped interval, with the two Q3 strips explicitly
budgeted, the original receiver field admits CLOCK's c-form with
the fixed core as its entire strong donor family. Capping its
donor mass at one costs the integrable core overshoot; its
normalized outside-core moment is identically zero. Two additional
selected-moment estimates complete CFR's field calculation.
The resulting value and induced receiver-derivative remainder
has integral at most a fixed multiple of $C_0E_g$ plus the strip
budget and a vanishing error. The full CLOCK forcing is small
for $E_g=2c_b$, small fixed $c_b,C_0$, and a small retained strip
budget. The actual entry reference, core cutoff and clock satisfy
CLOCK's passage hypotheses. The conclusion is conditional on
the other stopped inputs; no joint weak-passage bootstrap or
strip-entry theorem is claimed.
\end{abstract}

\section{Scope and precise stopped inputs}
All new labels have prefix \texttt{an03cfav1:}.
Use the original initialized flow and the notation of S2CP,
CFR and CLOCK:
\[
 \begin{gathered}
 P=\tfrac12(P_1+P_2),\quad P_1=N(e_1,q^2I_2),\quad
 P_2=N(\rho e_2,q^2I_2),\\
 \lambda=\rho^2,\quad Q=\lambda+q^2,\quad
 \rho\in[13/20,7/10],\quad\ell=\log(1/q).
 \end{gathered}
\]
Write $U=\sqrt m\,a_\theta$, $W=\sqrt m\,a_\psi$,
$f_S=\int_S W(U\cdot X)_+\dd\lambda_0$, and
\[
 \begin{gathered}
 E_S=\tfrac12\int_S(U_2^2+W_2^2),\qquad
 J_S=\tfrac12\int_S(U_1^2+W_1^2),\\
 \mu_S=E_S+J_S,\qquad
 c_S=Q^{-1}\E_{P_2}[(f_S)_2X_2].
 \end{gathered}
\]
The fixed partition is $\KK,\Gp,\Sigma,\BB$, with
$G=\Gp\cup\Sigma$ and $\mathcal D=\varnothing$.
Here $\KK$ is S2CP/TAILREP's fixed right-half core,
\begin{equation}
 \KK=\{\xi\le d_c\},\qquad
 d_c=c_*q^{\chi+kP/2},\quad \chi=1/20,\quad
 k=10(1-\lambda),\quad P=\Phi(\rho a),\quad a=\rho K(\rho a),
 \label{an03cfav1:cutoff}
\end{equation}
within the right-half ancestry, where
$K(z)=\phi(z)+z\Phi(z)$ and $\xi$ is actual-reference scaled
$\ell^1$ displacement at $t_0=t_h$.

On $[t_0,t_*]\subset[t_0,t_b]$ retain S2CP's stopped inputs:
\begin{equation}
 \begin{gathered}
 0\le c\le c_b<1,\qquad t_*-t_0\le L_T\ell,\qquad
 \int r_{\rm res}\le B,\quad r_{\rm res}=\sup_\theta|R_{11}|,\\
 J_G<q^2K_Y^2,\qquad E_{\Gp}<E_g,\qquad
 \|e_2\|_{P_1}+\|e_1\|_{P_2}\le M_{\rm res}q.
 \end{gathered}
 \label{an03cfav1:guards}
\end{equation}
Keep the core geometry $|\theta|,|\psi|\le Rq$,
$\mu_{\KK}\le M_K$, the actual same-field reference guard,
the S1 entries and the outside supremum of S2CP. Its fixed
receiver square is $[-L_c,L_c]^2$, with
$L_c=79/51+1$. The coarse fixed-chart c-form/remainder guard
used in S2CP is retained while its theorem is invoked. This
increment derives an explicit representation and a sharper
budget for that guard; it does not assume its own sharper
conclusion to prove S2CP's energy or mass statements.

Let $\mathcal S=\mathcal S_\pi\cup\mathcal S_{3\pi/2}$ be the
two fixed initial Q3 strips of GESG. The additional input is
exactly
\begin{equation}
 \sup_{[t_0,t_*]}\mu_{\mathcal S}/q\le b_{\mathcal S}(q)=o(1),
 \qquad \int_{t_0}^{t_*}\mu_{\mathcal S}/q\le F_{\mathcal S}.
 \label{an03cfav1:strips}
\end{equation}
For the field theorem $F_{\mathcal S}$ is any fixed finite
constant; the passage corollary requires its explicit smallness.
It is not proved here. No labels are reassigned.

The following are proved conclusions under these stopped inputs:
\begin{align}
 E_\Sigma&\le C_\Sigma(q^5+q^{\Gamma_\nu}H),&
 \mathcal U_\Sigma:=\int(E_\Sigma+K_Y\sqrt{E_\Sigma})&=o(1),
 \label{an03cfav1:S2inputs}\\
 \int\mu_{\BB}/q&\le F_{\mathcal S}
           +Cq^{\Delta_\nu}\ell+Cq^9\ell
           =F_{\mathcal S}+o(1),&
 \sup\mu_{\BB}/q&=o(1).\nonumber
\end{align}
They use S2CP Theorem S2, Corollary U and GESG's original-label
Q3 interior, with $\Gamma_\nu>3/10$, $\Delta_\nu>3/25$.
The outside decomposition is the strong complement, that
interior and the two strips. S2CP Proposition Bs also gives
\begin{equation}
 \mu_{\KK}(t_0)=\tfrac12-o(1),\quad
 w_+(t)+\int_{t_0}^t(\mu_{\KK}-1)_+\le d_q=o(1),
 \qquad w=\log\mu_{\KK}.
 \label{an03cfav1:overshootinput}
\end{equation}
These conclusions still retain core geometry. All errors are
uniform over stopped endpoints after the fixed constants are
chosen.

\section{Core-only donors and their mass normalization}
Let $\mu_s$ be the current radial measure of the fixed core,
pushed to $(\theta/q,\psi/q)$. Set
\begin{equation}
 \widetilde\mu_s=\mu_s\min(1,\mu_{\KK}^{-1}),\qquad
 \widetilde M=\min(\mu_{\KK},1),\qquad
 \widetilde Y=\int\eta\dd\widetilde\mu_s(\xi,\eta),\qquad
 \widetilde{\mathcal S}(x)
       =\int F(\rho x,\rho\xi)\dd\widetilde\mu_s,
 \label{an03cfav1:representation}
\end{equation}
where $F(z,w)=\phi(\min(z,w))+w\Phi(\min(z,w))$.
The reference is the original characteristic used by MOMC;
it is never reset or evolved in a collapsed field. Define
\[
 e=|(x_c,y_c)-(a,a)|_2,\qquad
 r_{\KK}=\mathop{\rm ess\,sup}_{b\in\KK}
 (|x_b-x_c|+|y_b-y_c|).
\]

\begin{lemma}[Permitted family choice and overshoot cost]
\label{an03cfav1:representationlemma}
CLOCK's cfield applies to the fixed strong family $\KK$ with
measure $\widetilde\mu_s$, once the field identity and remainder
below are proved. Its normalized tail moment is $T\equiv0$.
Replacing $\mu_s$ by $\widetilde\mu_s$ in the strong donor part
costs, in value plus induced $\ell^1$ receiver derivative,
\begin{equation}
 C_{L_c,R}(\mu_{\KK}-1)_+,\qquad
 \int C_{L_c,R}(\mu_{\KK}-1)_+\le C_{L_c,R}d_q=o(1).
 \label{an03cfav1:overshootcost}
\end{equation}
\end{lemma}
\begin{proof}
CLOCK cfield explicitly allows a fixed strong family with
positive instantaneous weights and requires no mass law or
constant normalized weights. Its proof is the pointwise donor
cancellation and the $G_M$ comparison of LOC; neither
differentiates those weights. Here $\widetilde M>0$, since
the entry core has positive mass and radial weights remain
positive on finite horizons. Uniform multiplication leaves the
normalized core law unchanged:
$\widetilde\mu_s/\widetilde M=\mu_s/\mu_{\KK}$.
The entire chosen family is the fixed core. Its complement in
that family is empty, so $T=0$ exactly. Physical labels outside
$\KK$ remain in the original field and must be paid separately.

Write $\delta=(\mu_{\KK}-1)_+$ and
$\nu_s=\mu_s/\mu_{\KK}$. The difference between unnormalized
and normalized strong fields, multiplied by two, is
\[
 \left(\delta x-\rho\delta\int F(\rho x,\rho\xi)\dd\nu_s,\quad
       \delta y-\delta\int\eta\dd\nu_s\right).
\]
On $|x|,|y|\le L_c$, $|\xi|,|\eta|\le R$, both the kernel
and its first receiver derivative
$\rho(\rho\xi-\rho x)_+\phi(\rho x)$ are uniformly bounded.
The displayed vector and its derivative matrix therefore
have the bound \eqref{an03cfav1:overshootcost}. Integrate
using \eqref{an03cfav1:overshootinput}. No derivative of the
time-dependent normalization factor is taken.
\end{proof}

\section{The additional selected moments needed for full assembly}
All norms in this section are
$\|h\|_{1,\infty}=\sup_{|x|\le L_c}(|h|+|h'|)$.
Under $P_2$ write $X=(qZ,V)$, with $Z$ standard Gaussian and
$V\sim N(\rho,q^2)$ independent. Let
$t_x=\tan(qx)/q$ and $I_x=\one_{\{Z+t_xV>0\}}$.
For a fixed cohort put
\[
 p_{1,S}=Q^{-1}\E[(f_S)_1V],\quad p_{2,S}=c_S,\qquad
 g_{i,S}=(f_S)_i-p_{i,S}V,
\]
\begin{equation}
 A_{1,S}=q^{-1}\E[g_{1,S}VI_x],\quad
 A_{2,S}=\E[g_{2,S}VI_x],\quad
 B_{2,S}=\E[g_{2,S}ZI_x].
 \label{an03cfav1:selecteddefs}
\end{equation}
$A_{2,S}$ is CFR's selected moment. The other two cannot be
discarded merely by global regression orthogonality.

\begin{lemma}[Full selected-moment budgets]
\label{an03cfav1:selected}
For an arbitrary cohort,
\begin{align}
 \|A_{1,S}\|_{1,\infty}
   &\le C_{L_c}(J_S+q^{-1}\sqrt{E_SJ_S}),\nonumber\\
 \|A_{2,S}\|_{1,\infty}+\|B_{2,S}\|_{1,\infty}
   &\le C_{L_c}(E_S+q\sqrt{E_SJ_S}).
 \label{an03cfav1:genericselected}
\end{align}
For $\Gp$, after the same fixed-chart decrease of $C_0$ as
in CFR, there is a nonnegative $\mathcal T$ such that
\begin{align}
 \|A_{1,\Gp}\|_{1,\infty}
 +\|A_{2,\Gp}\|_{1,\infty}+\|B_{2,\Gp}\|_{1,\infty}
 &\le C_{L_c}\{J_{\Gp}+q\sqrt{E_{\Gp}J_{\Gp}}
                         +\mu_{\Gp}\mathcal T\},\nonumber\\
 \int\mathcal T&\le C_{\rm RCET}(C_0+q^2L_T\ell).
 \label{an03cfav1:primaryselected}
\end{align}
An enlargement of the fixed Gaussian/RCET constants is understood.
In particular the primary integral is at most
$C_{L_c}C_{\rm RCET}C_0E_g+o(1)$, and the $\Sigma$ integral
is $C\mathcal U_\Sigma+o(1)=o(1)$.
\end{lemma}
\begin{proof}
Balance and label Cauchy--Schwarz give the pointwise bounds
\begin{align*}
 |(f_S)_1(qZ,V)|&\le2qJ_S|Z|+2\sqrt{E_SJ_S}|V|,\\
 |(f_S)_2(qZ,V)|&\le2q\sqrt{E_SJ_S}|Z|+2E_S|V|.
\end{align*}
These also control their regression coefficients. For any
continuous fixed function $h$ with these growth bounds,
boundary differentiation gives
\begin{align}
 \partial_x\E[hVI_x]
 &=t_x'\E_V[V^2h(-qt_xV,V)\phi(t_xV)],\nonumber\\
 \partial_x\E[hZI_x]
 &=-t_xt_x'\E_V[V^2h(-qt_xV,V)\phi(t_xV)].
 \label{an03cfav1:boundary}
\end{align}
The bounds $|t_x|\le2L_c$, $|t_x'|\le2$ and Gaussian moments
prove \eqref{an03cfav1:genericselected}, including the
$1/q$ in $A_{1,S}$. Polynomial domination justifies the
boundary formulas; the network output is continuous.

For the sharper primary bound write, exactly for $i=1,2$,
\[
 (f_{\Gp})_i=L_{i1}X_1+L_{i2}V+N_i,\quad
 L_{ij}=\int_{\Gp}W_iU_j,\quad
 N_i=\int_{\Gp}W_i(U\cdot X)_-.
\]
Here $|L_{11}|\le J_{\Gp}$ and
$|L_{21}|\le2\sqrt{E_{\Gp}J_{\Gp}}$.
Subtracting the regression removes $L_{i2}V$; the $X_1$
term has the stated $J_{\Gp}$ or $q\sqrt{E_{\Gp}J_{\Gp}}$
bound. We do not assume a sign for $W_1$.

For a unit input direction $u$ with $u_2>0$,
\[
 \E(u\cdot X)_-^2=q^2H_2(-\rho u_2/q)
       \le Cq^2e^{-c(u_2/q)^2}.
\]
Cauchy--Schwarz against $V$ or $Z$ and
$|W_i||U|\le m$ imply
\begin{equation}
 q^{-1}\E|N_iV|+q^{-1}\E|N_iZ|
 \le C\int_{\Gp}m e^{-c(u_2/q)^2}\dd\lambda_0.
 \label{an03cfav1:negativevalue}
\end{equation}
The same bound pays the subtracted regression. This explicit
factor $q$ in the negative-part moment is essential for
$A_{1,\Gp}$; an unscaled loss bound would not suffice.

RCET supplies $U_2>0$ and
$q|U_1|/U_2\le C(C_0e^{-b_*(t-t_0)/2}+q^2)$.
Choose $C_0$ so that $q|t_xU_1|\le U_2/2$ on this fixed
receiver square. On the boundary in \eqref{an03cfav1:boundary},
\[
 (U\cdot X)_-=[-(U_2-qt_xU_1)V]_+,
\]
which vanishes for $V\ge0$. For $V<0$ its contribution
to a first derivative is at most
$C_{L_c}\mu_{\Gp}q^{-1}\E[|V|^3\one_{\{V<0\}}]$,
where the factor $q^{-1}$ is only needed for $A_1$.
This is $C\mu_{\Gp}e^{-c'/q^2}$ after decreasing $c'$.
The regression part has the preceding value bound times
a bounded differentiated Gaussian moment.

RCET's growing margin therefore gives the same integrable
tail as CFR, with possibly smaller fixed exponential constant:
\[
 \mathcal T\le C\exp\!\left[
  -\frac{c}{q^2+C_*^2C_0^2e^{-b_*(t-t_0)}}\right]
                +Ce^{-c/q^2},\qquad b_*=Q(1-c_b)/2.
\]
Its integral is the second line of
\eqref{an03cfav1:primaryselected}, by the same Gaussian
inequality $e^{-c/v^2}\le C_cv$ and elementary exponential
integration used in CFR. Now use
$\mu_{\Gp}\le E_g+q^2K_Y^2$, $J_G\le q^2K_Y^2$ and
\eqref{an03cfav1:S2inputs}. This proves the final claims.
\end{proof}

\section{Full original-flow field assembly}
Let $O=\Gp\cup\Sigma\cup\BB$ and define the exact moments
\begin{equation}
 V_O=-\frac{\E_{P_2}[(f_O)_1X_2]}{qQ},\qquad
 K_O=\frac{\E_{P_1}[(f_O)_2X_1]}q,\qquad
 V=V_O,\quad K_w=(K_O)_+.
 \label{an03cfav1:VK}
\end{equation}
Use the normalized core in \eqref{an03cfav1:representation}.
The desired leading field is
\begin{align}
 2(F_c)_x&=-(1-\widetilde M)x+\rho\phi(\rho x)
       -\rho\widetilde{\mathcal S}(x)
       +\lambda[V+(1-c)y]\Phi(\rho x),\nonumber\\
 2(F_c)_y&=-(1-\widetilde M)y-\widetilde Y-K_w
                         +\rho(1-c)K(\rho x).
 \label{an03cfav1:cform}
\end{align}
Write $\|F\|_*=\sup_{[-L_c,L_c]^2}
(|F|_1+\|D_{x,y}F\|_{1\to1})$. This convention pays both
norms, including their sum in S2CP's fullnorm.

\begin{theorem}[Full fixed-chart remainder]
\label{an03cfav1:assembly}
Under the stopped inputs of Section 1 the exact scaled
receiver velocity satisfies
$F_{\rm orig}=F_c+\mathcal R$, with
$\eps_F(t)=\|\mathcal R(t)\|_*$ and
\begin{equation}
 \int_{t_0}^{t_*}\eps_F
 \le C_F C_{L_c}C_{\rm RCET}C_0E_g
                    +C_{\rm out}F_{\mathcal S}+o(1).
 \label{an03cfav1:fullbudget}
\end{equation}
$C_F,C_{L_c},C_{\rm out}$ depend only on the fixed chart,
core radius/mass and model/Gaussian constants.
Residual-guard dependence is allowed in $C_{\rm RCET}$ and
in coefficients of vanishing scales. No selected residual
is removed by regression orthogonality.
\end{theorem}
\begin{proof}
We give a componentwise assembly, including the moments
not present in CFR's second-output selected estimate.
The exact formulas are
\[
 2q(F_{\rm orig})_x
 =\sum_{j=1}^2\E_{P_j}[(b_y\cdot(X-f))
                    (a_x^\perp\cdot X)I_x],\qquad
 2q(F_{\rm orig})_y
 =\sum_{j=1}^2\E_{P_j}[(b_y^\perp\cdot(X-f))(a_x\cdot X)I_x],
\]
where $a_x=(\cos qx,\sin qx)$ and similarly for $b_y$.
The gates here are interpreted under the relevant cluster.
The mixture factor is accounted for by the left factor two.

First retain only target and core output. MOMC's reduction
proof is linear in the donor measure; replacing its bound
$M\le1$ by the retained $M_K$ simply changes its error to
$C_{L_c,R}q^2(1+M_K)$ (use $\max(L_c,R)$ for its two compact
radii). It gives
\[
 \begin{split}
 2(F_{\rm core})_x&=-(1-\mu_{\KK})x+\rho\phi(\rho x)
              -\rho\mathcal S(x)+\lambda y\Phi(\rho x),\\
 2(F_{\rm core})_y&=-(1-\mu_{\KK})y-Y+\rho K(\rho x),
 \end{split}
\]
with this error in value and first receiver derivatives.
In particular the full order-one core overlap is kept.
The core's second output on $P_2$ is
$O_{C^1}(q^2\mu_{\KK})$ by the same kernel proof and CFR
moment; $|c_{\KK}|\le C_Rq^2\mu_{\KK}$.

For $G=\Gp\cup\Sigma$, first consider $P_1$.
Receiver gates can be removed in value and first derivative
at cost $C\mu_G q^{-N}e^{-c/q^2}$ for a fixed $N$.
This follows from the bounded-chart boundary argument in
MOMC: the output bound is $\mu_G|X|$, and its boundary
moments have a Gaussian factor $e^{-c/q^2}$.
The exact Stein identities and elementary output bound give
\begin{align}
 \E_{P_1}[(f_G)_iX_2]
 &=q^2\int_GW_iU_2\,\mathbb P_{P_1}(U\cdot X>0)\dd\lambda_0,
 \nonumber\\
 |\E_{P_1}[(f_G)_1X_1]|
 &\le C(J_G+q\sqrt{E_GJ_G}).
 \label{an03cfav1:P1stein}
\end{align}
Consequently the leading $P_1$ non-core contributions to
$2F_x,2F_y$ are $0,-K_G$, respectively, with error at most
\begin{equation}
 C_{L_c}\{J_G+q\sqrt{E_GJ_G}+q^2E_G+q^2|K_G|\}
                  +C\mu_G q^{-N}e^{-c/q^2}
 \label{an03cfav1:P1error}
\end{equation}
in value and first derivatives. These estimates follow by
expanding the displayed exact trigonometric factors; all
their discarded coefficients are $O_{C^1}(q^2)$ after the
leading $1/q$ has been accounted for in $K_G$.

On $P_2$, write
$T_2(x)=\E[V^2I_x]$ and $T_1(x)=\E[VZI_x]$.
The Gaussian identities used in MOMC give
\begin{equation}
 T_2=\lambda\Phi(\rho x)+O_{C^1}(q^2),\qquad
 T_1=\rho\phi(\rho x)+O_{C^1}(q^2).
 \label{an03cfav1:gaussian}
\end{equation}
The first-order trigonometric expansion of the non-core
subtraction is, component by component,
\begin{align}
 2(F_G^{P_2})_x
 &=V_GT_2-yc_GT_2-A_{1,G}-yA_{2,G}
                              +\mathcal E_x,\nonumber\\
 2(F_G^{P_2})_y
 &=-c_G(T_1+xT_2)-B_{2,G}-xA_{2,G}
                              +\mathcal E_y.
 \label{an03cfav1:P2assembly}
\end{align}
Here $V_G=-p_{1,G}/q$. The error has norm at most
$C_{L_c}q^2(E_G+J_G+q^{-1}\sqrt{E_GJ_G})$.
For clarity, this follows directly from
$a_x^\perp\cdot X=V+O_{C^1}(q^2)(|V|+|Z|)$,
$a_x\cdot X=q(Z+xV)+O_{C^1}(q^3)(|V|+|Z|)$,
$b_y=(1,qy)+O_{C^1}(q^2,q^3)$ and
$b_y^\perp=(-qy,1)+O_{C^1}(q^3,q^2)$.
Use the two output bounds in Lemma \ref{an03cfav1:selected}
also on the receiver boundary to control differentiated
expectations. Thus no unproved trace estimate is used.

By \eqref{an03cfav1:gaussian},
$T_1+xT_2=\rho K(\rho x)+O_{C^1}(q^2)$.
Lemma \ref{an03cfav1:selected} pays every selected term
in \eqref{an03cfav1:P2assembly}.
The residual $X_2$ coefficient is $c_G$ at this stage,
not silently the full $c$.

For $\BB$ apply PROFILE Proposition 4.1 to its exact
two-component velocity and first receiver derivative.
Its cost is $C\mu_{\BB}/q$. Replacing that output by its
formal contributions $V_{\BB},-K_{\BB},-c_{\BB}$ costs
the same order, because Gaussian moments give
$|V_{\BB}|+|K_{\BB}|+|c_{\BB}|\le C\mu_{\BB}/q$
and the coefficient functions on the fixed square have
bounded $C^1$ norms. Now $V_G+V_{\BB}=V_O$,
$K_G+K_{\BB}=K_O$, and
$c_G+c_{\BB}=c-c_{\KK}$. Replacing this last coefficient
by $c$ costs $Cq^2M_K$. Replace $K_O$ by $(K_O)_+$ and
put the resulting $(K_O)_-/2$ into the second remainder
component. CFR's negative-part lemma bounds its integral
by $C\mathcal U_\Sigma+C\int\mu_{\BB}/q+o(1)$.
Finally normalize the core by Lemma
\ref{an03cfav1:representationlemma}.

For completeness, a single pointwise majorant for all these
costs, with a fixed enlargement of $C$, is
\begin{align}
 \eps_F\le C\big[&
 q^2(1+M_K+M_{\rm res}+\mu_G+|V_O|+|K_O|)
 +J_G+q\sqrt{E_GJ_G}
 +\mu_{\Gp}\mathcal T\nonumber\\
 &+E_\Sigma+q^{-1}\sqrt{E_\Sigma J_\Sigma}
 +\mu_{\BB}/q+(\mu_{\KK}-1)_+
 +\mu_G q^{-N}e^{-c/q^2}\big].
 \label{an03cfav1:majorant}
\end{align}
The exponentially small primary part of $(K_O)_-$ is
included in the last term. The added $q^2M_{\rm res}$ is
a harmless enlargement: if matrix-entry estimates are used
for remaining cross-cluster terms, the extra chart factors
give this order. Bare $O(M_{\rm res}q)$ cross entries divided
by $q$ are \emph{not} discarded; their order-one moments
are $V_O,K_O$ and the selected terms already retained.

Integrate \eqref{an03cfav1:majorant}.
The primary selected term costs
$C C_{L_c}C_{\rm RCET}C_0E_g+o(1)$.
$J_G$ and $q\sqrt{E_GJ_G}$ cost
$O(q^2\ell(K_Y^2+K_Y\sqrt{E_g+1}))$.
The secondary terms cost $C\mathcal U_\Sigma$.
S2CP's candidate moment bound, or its instantaneous
cross-output bound and outside supremum, pays the $q^2$
multiple of $|V_O|+|K_O|$ by $o(1)$.
The remaining terms are paid by
\eqref{an03cfav1:S2inputs} and
\eqref{an03cfav1:overshootinput}. This proves
\eqref{an03cfav1:fullbudget}.
\end{proof}

\paragraph{Integrated term ledger.}
All entries below include value and induced receiver derivative;
constants are fixed before $q$.
\begin{longtable}{>{\raggedright\arraybackslash}p{.42\textwidth}
 >{\raggedright\arraybackslash}p{.50\textwidth}}
\textbf{Term}&\textbf{Integrated price and source}\\\hline
Target/core finite-noise kernels&
$C_{L_c,R}q^2(1+M_K)L_T\ell$; MOMC's compact-kernel proof.\\
Core second output and $c_{\KK}$&
$Cq^2M_KL_T\ell$; CFR moment and MOMC.\\
Core overshoot normalization&
$C_{L_c,R}d_q=o(1)$; S2CP's positive-part identity.\\
$G_+$ selected residuals of both components&
$C C_{L_c}C_{\rm RCET}C_0E_g+o(1)$;
CFR and Lemma \ref{an03cfav1:selected}.\\
$P_1$ transverse moments, $q^2$ corrections and chart gate tails&
$Cq^2\ell(1+M_K+M_{\rm res}+K_Y^2+K_Y\sqrt{E_g+1})
 +o(1)$; \eqref{an03cfav1:P1stein}--\eqref{an03cfav1:P2assembly}.\\
$\Sigma$ selected moments and its negative $K_O$ part&
$C\mathcal U_\Sigma+o(1)$; S2CP U, CFR negative and
\eqref{an03cfav1:genericselected}.\\
Strong complement&
$Cq^{\Delta_\nu}\ell$; S2CP compmass and PROFILE.\\
Original Q3 interior&
$Cq^9\ell$; GESG and PROFILE.\\
Both Q3 strips, including their moment substitutions&
$C_{\rm out}F_{\mathcal S}$; retained strip budget and PROFILE.\\
Primary negative $K_O$ part&
Gaussian $o(1)$; CFR negative.\\
\end{longtable}

\section{Uniform CLOCK constants and its forcing budget}
The output of \eqref{an03cfav1:representation} is an actual
cfield representation of the common original receiver field,
with $0<\widetilde M\le1$, $K_w\ge0$ and $T=0$.
Let $\alpha=\lambda P/2$.
LOC's coherent matrix obeys
\[
 J_M=J_1-\operatorname{diag}((1-M)(1/2-\alpha),0),\qquad
 z^TJ_Mz\le-\gamma|z|^2,\quad
 \gamma=\frac{\alpha(1/2-\alpha)}{\alpha+1/2}.
\]
The proved box bounds used in MOMC give
$\gamma\ge\gamma_*:=1183/20800>0$.
Choose a fixed enlargement $C_{\rm pass}$ of the uniform
LOC constant $724/51$ and the remainder norm-conversion
constants, and set
\begin{equation}
 g=\gamma_*/2,\qquad
 e_*=\min\{1/4,\gamma_*/(2C_{\rm pass})\},\qquad
 K_{\rm pass}=C_{\rm pass}(1+C_{\rm pass}/g).
 \label{an03cfav1:passconstants}
\end{equation}
These depend only on the model box and fixed neighborhood,
not on $c_b,C_0$ or residual guards. If $e<e_*$ and
$r_{\KK}<1$, all core/reference points and joining segments
lie inside $[-L_c,L_c]^2$, since $a<28/51$.
CLOCK cfield and LOC's donor cancellation therefore give
\begin{equation}
 \begin{split}
 D^+e&\le-ge+C_{\rm pass}(r_{\KK}+f_*),\\
 D^+r_{\KK}&\le
 [\alpha(1-c)+C_{\rm pass}(e+r_{\KK}+f_*)]r_{\KK},\\
 f_*&=c+|V_O|+(K_O)_++\eps_F .
 \end{split}
 \label{an03cfav1:localineq}
\end{equation}
Time-dependent donor weights cause no additional derivative
term: these are instantaneous same-field comparisons, not
equations for a weighted barycenter.

\begin{proposition}[Forcing bound and its constant order]
\label{an03cfav1:forcing}
Under the same stopped hypotheses,
\begin{align}
 \int f_*
 \le{}&\frac{10e}{9}\frac{E_g}{Q(1-c_b)}
 +\frac{C_V}{Q(1-c_b)}(E_g+K_Y\sqrt{E_g})\nonumber\\
 &+C_F C_{L_c}C_{\rm RCET}C_0E_g
 +C_{\mathcal S}F_{\mathcal S}+o(1).
 \label{an03cfav1:forcingbudget}
\end{align}
For $E_g=2c_b$ this is
$\kappa(c_b,C_0)+C_{\mathcal S}F_{\mathcal S}+o(1)$,
where $\kappa\to0$ as $c_b,C_0\to0$, with uniform residual
guard constants as specified below.
\end{proposition}
\begin{proof}
Use S2CP's c budget, retaining $1/Q$, and its VKbudget.
The latter is
$C_V(E_g+K_Y\sqrt{E_g})/b_0+C\int\mu_{\BB}/q+o(1)$,
$b_0=Q(1-c_b)$. Add \eqref{an03cfav1:fullbudget} and
$T=0$, then use \eqref{an03cfav1:S2inputs}. In particular
one may take
\[
 \kappa(c_b,C_0)=
 \frac{20e\,c_b/9+C_V(2c_b+K_Y\sqrt{2c_b})}
      {(169/400)(1-c_b)}
       +2C_F C_{L_c}C_{\rm RCET}C_0c_b.
\]
The asserted limiting property follows once the displayed
constants are fixed uniformly over small $c_b$.

Here is an explicit such choice. Fix $\bar c=1/4$, the
model/core budgets and $B_s>\log2$ first. In LCRC's constant
construction replace its old energy budget $2+2c_b$ by
$2+2\bar c$, its $h_*=(9/8)c_b+1$ by
$\bar h_*=(9/8)\bar c+1$, and its lower clock rate by
$\underline b=(169/400)(1-\bar c)$.
Set $B=C_rB_s+2$ and use these larger finite numbers in
LCRC's displayed $E_*,K_E,K_{1/2},A_*,K_Y,M_{\rm res}$
choices. All are independent of the subsequently selected
$0<c_b\le\bar c$. The smaller energy guard $E_g=2c_b$
has the pointcap slack checked in S2CP. RCET's constants
and admissible small $C_0$ can likewise use the fixed
positive lower rate $\underline b/2$ throughout this range.
There is no singular $1/c_b$ in these choices.
\end{proof}

For clarity, the small-budget selection is not a claim that
the strip estimate is automatic. Set
$P_{\rm allow}=e_*/(4C_{\rm pass})$ and choose a fixed
coarse field allowance $0<\overline F_F<\infty$ before
the S2 distance constant. Choose $c_b\le\bar c$ so that
the first two terms of \eqref{an03cfav1:forcingbudget}
are at most $P_{\rm allow}/4$.
Choose $C_0$ within RCET's uniform admissibility and the
receiver-boundary margin, with
\[
 C_F C_{L_c}C_{\rm RCET}C_0E_g
       \le\min(P_{\rm allow}/4,\overline F_F/4).
\]
Require the retained strip budget to satisfy
\begin{equation}
 C_{\mathcal S}F_{\mathcal S}
       \le\min(P_{\rm allow}/4,\overline F_F/4),
 \label{an03cfav1:stripsmall}
\end{equation}
where $C_{\mathcal S}$ is enlarged to dominate the strip
constant in the field bound. Choose $q$ last so all remaining
errors are at most $\min(P_{\rm allow}/4,\overline F_F/4)$.
Then $\int f_*\le P_{\rm allow}$ and
$\int\eps_F\le3\overline F_F/4$.

A consistent complete order is
\begin{equation}
 \begin{gathered}
 (\text{model},R,M_K,B_s,L_c,\bar c,C_{\rm pass},g,e_*)
 \ \longrightarrow\ (B,K_Y,M_{\rm res})\\
 \longrightarrow\ (\overline F_F,c_b,E_g=2c_b,C_D,A_0,C_m)
 \ \longrightarrow\ C_0\\
 \longrightarrow\ (K(C_0),Z_\Sigma,\text{entry constants},C_\Sigma)
 \ \longrightarrow\ q .
 \end{gathered}
 \label{an03cfav1:order}
\end{equation}
The required fixed small upper bound on $F_{\mathcal S}$ is
chosen before $q$. A given larger strip budget still yields
\eqref{an03cfav1:fullbudget} and
\eqref{an03cfav1:forcingbudget}, but not the small-forcing
passage conclusion. No guard-independent chart or core
constant is enlarged in response to $C_\Sigma$ or $C_0$.

\section{Entry audit and conditional passage}
\begin{corollary}[CLOCK passage with the actual reference and clock]
\label{an03cfav1:passage}
Retain the non-passage stopped inputs of Section 1 and the
strip bounds with \eqref{an03cfav1:stripsmall}; choose constants
as above. On any such interval ending at $t_*\le t_b$, for
sufficiently small $q$, CLOCK's passage theorem applies with
clock cohort $\mathcal C_q$, strong family/core $\KK$ and
the representation \eqref{an03cfav1:representation}.
It gives
\begin{equation}
 e(t)\le3e_*/4,\qquad r_{\KK}(t)\le2\mathcal A,\qquad
 \int_{t_0}^{t_*}r_{\KK}\le2\mathcal A/a_b,
 \quad \mathcal A=O(q^\chi),\quad a_b=\alpha(1-c_b).
 \label{an03cfav1:passresult}
\end{equation}
This is a conditional passage improvement on the common
stopped trajectory. It does not continue the other guards
to $t_b$, prove the strip bounds, or execute the joint
weak-passage bootstrap.
\end{corollary}
\begin{proof}
MOMC proves for the actual reference at $t_0=t_h$ that
$e_0=o(1)$, hence $e_0\le e_*/4$.
By definition of the fixed core, $r_0\le d_c$.
The clock is the original $\eta=1/8$ cohort, not all of
$\Gp$: LCRC localclock/BULK give positivity and
\[
 \int_{t_0}^{t_*}\eps_{\mathcal C_q}
       \ge-D_-,\qquad D_-=O(q^{1/8}+q^2\ell)=o(1).
 \]
S2CP's energy conclusion and the LCRC pointcap lemma
give $H(t_*)\le(9/8)c_b+o(1)$ on the stopped interval,
so take the fixed $H_b=(5/4)c_b$ for small $q$.
The S1 seed is
$H_0\asymp A_q^{-\lambda}q^k$, with
$0<c_A\le A_q\le C_A$.
These are conclusions under the current guards, not
assumed endpoint caps or independent clock premises.

Use CLOCK's exact factor
\[
 L_b=\exp(C_{\rm pass}e_0/g+K_{\rm pass}F_b+\alpha D_-/Q)
                      (H_b/H_0)^{\alpha/Q},
 \qquad F_b=P_{\rm allow}.
 \]
With the cutoff \eqref{an03cfav1:cutoff},
\begin{align}
 r_0L_b
 &\le Cc_*H_b^{\alpha/Q}A_q^{\lambda\alpha/Q}
        q^{\chi+kP/2-k\alpha/Q}\nonumber\\
 &=Cc_*H_b^{\alpha/Q}A_q^{\lambda\alpha/Q}
        q^{\chi+kPq^2/(2Q)}
 \le C' q^\chi.
 \label{an03cfav1:handoff}
\end{align}
The constant includes the fixed seed comparison and bounded
exponential factor. The positive $q^2$ exponent is retained;
no false equality $\alpha/Q=P/2$ is used.
Since $\alpha\ge169/1600$ and $c_b\le\bar c$,
$a_b$ has a fixed positive lower bound. Consequently
\[
 r_0L_b\le
 \min\{1/4,\ a_b/(2K_{\rm pass}),\
                        ge_*/(8C_{\rm pass})\}
 \]
for small $q$. The forcing budget was proved in
Proposition \ref{an03cfav1:forcing}. Finally the mean and
shape inequalities are \eqref{an03cfav1:localineq}, obtained
from the now-proved actual c-form, not from a renamed weak
mass. All hypotheses of CLOCK's passage theorem therefore
hold at the endpoint $t_*$. Apply that theorem to get
\eqref{an03cfav1:passresult}.
\end{proof}

\paragraph{Boundary of the result.}
The four requested tasks close as implications under the
specified stopped inputs and retained strip budgets.
The old coarse field guard was used only to invoke S2CP,
whose energy/mass conclusions then permit its explicit
refinement. The resulting strict field and passage
improvements are suitable for a later joint assembly;
they are not that assembly. No new weak-entry contract,
strip classification, selection theorem or consolidation
is claimed. The companion review records source provenance,
constant dependencies and each retained input.
\end{document}
